% ---- pg-384 （书页 372） ----

$x_1$: Family (家庭中人口数)

$x_2$: DistRd (到最近大路的距离(km))

$x_3$: Cotton (2000 年种植棉花的公顷数)

$x_4$: Maize (2000 年种植玉米的公顷数)

$x_5$: Sorg (2000 年种植高粱的公顷数)

$x_6$: Millet (2000 年种植黍子的公顷数)

$x_7$: Bull (饲养的骡牛或其他骡家畜的总头数)

$x_8$: Cattle (奶牛总数)

$x_9$: Goats (山羊总数)

数据见\hyperref[tab:8.7]{表 8.7} (全部数据见网页 www.prenhall.com/statistics).

\begin{mvtable}
\phantomsection\label{tab:8.7}\mvtabcap{表 8.7\quad 马里家庭农场数据}
\begin{tabular}{rrrrrrrrr}
\toprule
Family & DistRd & Cotton & Maize & Sorg & Millet & Bull & Cattle & Goats\\
\midrule
12&80&1.5&1.00&3.0&0.25&2&0&1\\
54&8&6.0&4.00&0&1.00&6&32&5\\
11&13&0.5&1.00&0&0&0&0&0\\
21&13&2.0&2.50&1.0&0&1&0&5\\
61&30&3.0&5.00&0&0&4&21&0\\
20&70&0&2.00&3.0&0&2&0&3\\
29&35&1.5&2.00&0&0&0&0&0\\
29&35&2.0&3.00&2&0&0&0&0\\
57&9&5.0&5.00&0&0&4&5&2\\
23&33&2.0&2.00&1.0&0&2&1&7\\
\multicolumn{9}{c}{$\vdots$}\\
20&0&1.5&1.00&3.0&0&1&6&0\\
27&41&1.1&0.25&1.5&1.50&0&3&1\\
18&500&2.0&1.00&1.5&0.50&1&0&0\\
30&19&2.0&2.00&4.0&1.00&2&0&5\\
77&18&8.0&4.00&6.0&4.00&6&8&6\\
21&500&5.0&1.00&3.0&4.00&1&0&5\\
13&100&0.5&0.50&0&1.00&0&0&4\\
24&100&2.0&3.00&0&0.50&3&14&10\\
29&90&2.0&1.50&1.5&1.50&2&0&2\\
57&90&10.0&7.00&0&1.50&7&8&7\\
\bottomrule
\end{tabular}

{\footnotesize 资料来源: Data courtesy of Jay Angerer.}
\end{mvtable}

(a) 构造 Family-DistRd 和 DistRd-Cattle 的二维散点图. 从数据集中去除任何明显的离群值.

(b) 对相关阵 $\mathbf{R}$ 作主成分分析. 确定适当数量的主成分以有效地综合变差. 用解释的变差比例和碎石图来帮助你确定维数.

(c) 解释前五个主成分. 你是否能识别类似于“农场规模”那样的成分? 或者类似于“山羊与到大路距离”那样的成分?
